\par\SourceLeaf{497}%
\begingroup\fontsize{10.5}{13}\selectfont\interlinepenalty=10000
\par\noindent\hangindent=1em Abrisse der Messungen im Bremischen 1839. 1843. 1844 . . . . . . . . . . .\hfill\penalty0\hbox{}\nobreak\hfill\mbox{Seite 478}\par
\par\noindent\hangindent=1em Abrisse aus der Vereinigung der Messungen vom Jahre 1843. 1844 . . . . . . . .\hfill\penalty0\hbox{}\nobreak\hfill\mbox{— 480}\par
\par\noindent\hangindent=1em Bericht über die Resultate der trigonometrischen Messungen . . . . .\hfill\penalty0\hbox{}\nobreak\hfill\mbox{1848 März 15 — 481}\par
\par\noindent\hangindent=1em Entwurf zur Gradmessung . . . . . . . . . . . . . . . . .\hfill\penalty0\hbox{}\nobreak\hfill\mbox{1819 Mai 30 — 484}\par
\par\noindent\hangindent=1em Über die Fortsetzung der Dänischen Gradmessung durch Hannover . . .\hfill\penalty0\hbox{}\nobreak\hfill\mbox{1820 Aug. 11 — 485}\par
\par\noindent\hangindent=1em Über \textsc{Gauss} Gegenwart an einem Theil der Operationen der Dänischen Gradmessung . . . . . . . . . . . . . . . . . . . . .\hfill\penalty0\hbox{}\nobreak\hfill\mbox{1820 Febr. 27 — 486}\par
\par\noindent\hangindent=1em Promemoria betreffend die Hannoversche Gradmessung . . . . . . .\hfill\penalty0\hbox{}\nobreak\hfill\mbox{1821 März 20 — 487}\par
\endgroup
\par\bigskip\Needspace{7\baselineskip}\begin{center}GÖTTINGEN,\par\medskip GEDRUCKT IN DER DIETERICHSCHEN UNIVERSITÄTS-DRUCKEREI\par\medskip\textsc{W. Fr. Kaestner.}\end{center}
